% problem_id: S047-b21-Q4
% source_id: S047/S048 (numdam/ASNSP_2013_5_12_4_1001_0.pdf)
% source_file: ASNSP_2013_5_12_4_1001_0.pdf
% statement_locator: printed p. printed p. 1009 (PDF p. 9): the subsection heading "3.2. Key estimates", the framing sentence, the heading "Lemma 3.3 (Magic constants).", items (1) and (2) with the boxed criterion (3.2), and the closing words "holds true." — the whole statement lies on this page.
% proof_locator: printed p. printed pp. 1009-1010 (PDF pp. 9-10): "Proof." begins at the foot of printed p. 1009 (immediately after (3.2)) and the closing square is on printed p. 1010; the page break falls between the display "-1 <= P_n(x) <= P_n(0)+x = 1+x <= 2" (p. 1009) and "for every x in [0,1]. In particular, ..." (p. 1010). Pages read as images for this candidate: PDF 5, 9, 10, 11, 12 (= printed 1005, 1009, 1010, 1011, 1012).
% representation: PDF; transcribed from rendered page images
% collected_by: carrier rule
% review_status: H2 by 2/2 independent reviewers (the miner claimed H3); 2/2 忠实
% status: complete (statement + author proof verbatim + reason + locators)
%
% =====================================================================
% Candidate: Q4
% Paper: Marina Ghisi, Massimo Gobbino, "Higher order Glaeser inequalities
%        and optimal regularity of roots of real functions",
%        Ann. Sc. Norm. Super. Pisa Cl. Sci. (5) XII (2013), 1001-1021 (21 pp.).
%        numdam: https://www.numdam.org/item/ASNSP_2013_5_12_4_1001_0.pdf
% Source file: /disks/sata1/yupeng/human-proof-corpus/source-cache/registry-sources/numdam/ASNSP_2013_5_12_4_1001_0.pdf
% Rendered with: pdftoppm -f <p> -l <p> -r 150 -png   (pdftotext used ONLY for locating;
%   its text layer mangles the signs: on printed p. 1009 it prints "P(0) = 1 and P'(0) = 1"
%   for "P(0) = 1 and P'(0) = -1", "Pn(x)  1" for "P_n(x) >= -1", and "A  "k" for "A >= e_k")
% Page mapping: PDF page k == printed page 1000 + k.
%
% COLLECTED UNIT (see Q4.json): Lemma 3.3 (Magic constants), printed pp. 1009-1010
% (= PDF pp. 9-10): statement items (1)-(2) + boxed criterion (3.2) + the complete author proof.
%
% WHY THIS ONE (the carrier, not the signboard):
%   * Theorem 2.1 (printed p. 1005), the first flagship, is proved on printed p. 1012 by ONE
%     paragraph that hands everything to Proposition 3.4 and then lets d -> +infinity
%     ==> reduction-grade (H2, Context A + Context B below).
%   * Proposition 3.4 is ALREADY COLLECTED, and its own record in this corpus downgraded it to
%     H2 for exactly one reason: both of its quantitative kernels are quoted, not proved there --
%     the scale delta_k comes from Lemma 3.3 ("Let delta_k be the constant of Lemma 3.3",
%     printed p. 1011) and the threshold epsilon_k from criterion (3.2) ("Therefore from (3.2) it
%     follows that A >= epsilon_k", printed p. 1012).  Proposition 3.5 (also already collected)
%     sits further down the same chain.  Lemma 3.3 is the terminal node of that deferral chain:
%     it quotes nothing beyond finite-dimensionality / Vandermonde, and it produces both constants.
%
% Contexts A-C below are quoted verbatim ONLY so that this chain can be checked on the page;
% they are NOT part of the collected unit.
%
% Verbatim transcription. No rewriting, no completion. Original typos are kept and flagged.
% Illegible spots are flagged as %% [ILLEGIBLE: ...].
% =====================================================================

\documentclass[11pt]{article}
\usepackage{amsmath,amssymb,amsthm}
\usepackage[margin=1in]{geometry}
\newtheorem{lemma}{Lemma}
\newcommand{\Hold}{\mathrm{H\ddot{o}ld}}
\begin{document}

\section*{Q4 --- Transcription (verbatim)}

% ---------------------------------------------------------------------
% CONTEXT A (verbatim, printed p. 1005): the flagship whose proof is a reduction.
% ---------------------------------------------------------------------
\subsection*{Context A --- Theorem 2.1 (printed p.~1005, PDF p.~5), the flagship}

\textbf{Theorem 2.1 (Higher order Glaeser inequalities).} {\itshape
Let $k$ be a positive integer, let
$\alpha \in (0,1]$, and let $v \in C^k(\mathbb{R})$ be a function such that $v(x)$ and $v'(x)$ do
not change their sign in $\mathbb{R}$ (namely either $v(x) \ge 0$ or $v(x) \le 0$ for every
$x \in \mathbb{R}$, and either $v'(x) \ge 0$ or $v'(x) \le 0$ for every $x \in \mathbb{R}$).

Let us assume that the $k$-th derivative $v^{(k)}(x)$ is $\alpha$-H\"older continuous in
$\mathbb{R}$.

Then there exists a constant $C(k)$, depending only upon $k$, such that
}
\begin{equation}
|v'(x)|^{k+\alpha} \;\le\; C(k)\cdot|v(x)|^{k+\alpha-1}\cdot \Hold_\alpha\!\left(v^{(k)},\mathbb{R}\right)
\qquad \forall x \in \mathbb{R}. \tag{2.3}
\end{equation}

% ---------------------------------------------------------------------
% CONTEXT B (verbatim, printed p. 1012): the ENTIRE proof of the flagship.
% ---------------------------------------------------------------------
\subsection*{Context B --- the entire proof of Theorem 2.1 (printed p.~1012, PDF p.~12)}

\emph{Proof of Theorem 2.1.} Let us fix any point $x_0 \in \mathbb{R}$. Then the assumptions of
Proposition 3.4 are satisfied for every $d>0$. The conclusion follows by remarking that
\[
\Hold_\alpha\!\left(v^{(k)},(x_0-d,\,x_0+d)\right) \;\le\; \Hold_\alpha\!\left(v^{(k)},\mathbb{R}\right),
\]
and finally letting $d \to +\infty$. \hfill$\square$

% ---------------------------------------------------------------------
% CONTEXT C (verbatim, printed pp. 1011-1012): where the already collected Proposition 3.4
% consumes (rather than produces) the two kernels that Lemma 3.3 produces.
% ---------------------------------------------------------------------
\subsection*{Context C --- where Proposition 3.4 (already collected) quotes Lemma 3.3}

%% [Editorial note, not transcription: all four quotations below are sentences of the proof of
%%  Proposition 3.4 (printed p. 1011 and printed p. 1012).  They are reproduced for the sole
%%  purpose of recording, on the page, that delta_k and (3.2)/epsilon_k are consumed there.]

\noindent (printed p.~1011) Let $\delta_k$ be the constant of Lemma 3.3. We distinguish two cases.

\noindent (printed p.~1011) \dots which implies (3.5) in this first case, provided that
$C(k) \ge \delta_k^{\,k}$.

\noindent (printed p.~1012) \dots we have that $P(0)=1$, $P'(0)=-1$, and the assumptions in the
left-hand side of (3.2) are satisfied.

\noindent (printed p.~1012) Therefore from (3.2) it follows that $A \ge \varepsilon_k$, which
implies (3.5) also in this second case, provided that $C(k) \ge (\varepsilon_k k!)^{-1}$.
\hfill$\square$

% =====================================================================
% THE COLLECTED UNIT
% =====================================================================
\subsection*{Collected unit --- Lemma 3.3 (Magic constants), printed pp.~1009--1010 (PDF pp.~9--10)}

%% --- page 1009 (PDF page 9; running head: OPTIMAL REGULARITY OF ROOTS) ---
%% [The part of printed p. 1009 above the heading 3.2 is the closing paragraph of the proof of
%%  Lemma 3.2 (its last display is the integration-by-parts computation ending in "which
%%  completes the proof."); it is not part of this unit, so the transcription starts at the
%%  heading 3.2.]

\subsection*{3.2.\quad Key estimates}

The following surprising property of polynomials is the key tool in the proof of our higher order
Glaeser inequalities. We stress that the ``magic constants'' $\delta_k$ and $\varepsilon_k$ do not
depend on $P(x)$.

\renewcommand{\thelemma}{3.3}
\begin{lemma}[Magic constants]
\emph{For every positive integer $k$ there exist positive constants $\delta_k$ and $\varepsilon_k$,
depending only upon $k$, with the following properties.}

(1) \emph{For every polynomial $P(x)$, with degree less than or equal to $k$, such that $P(0)=1$
and $P'(0)=-1$, there exists $x \in [0,\delta_k]$ such that either $P(x) \le -1$ or
$P'(x) \ge 1$.}

(2) \emph{For every real number $A$, every $\alpha \in (0,1]$, and every polynomial $P(x)$, with
degree less than or equal to $k$, such that $P(0)=1$ and $P'(0)=-1$, the following implication}
\begin{equation}
\left.
\begin{array}{ll}
P(x)+Ax^{k+\alpha} \ge 0 & \forall x \in [0,\delta_k]\\[3pt]
P'(x)-Akx^{k+\alpha-1} \le 0 & \forall x \in [0,\delta_k]
\end{array}
\right\}
\;\Longrightarrow\;
\boxed{\;A \ge \varepsilon_k\;}
\tag{3.2}
\end{equation}
\emph{holds true.}
\end{lemma}

%% [Editorial note, not transcription: in print the two lines of the hypothesis of (3.2) are
%%  enclosed in one box and the conclusion "A >= epsilon_k" in a second box, with the double
%%  arrow printed between the two boxes.]

\subsection*{Proof (authors' text)}

Let us assume that statement (1) is false for some $k$. Then there exists a sequence of polynomials
$P_n(x)$, with degree less than or equal to $k$, such that $P_n(0)=1$, $P_n'(0)=-1$, and
\begin{equation}
P_n(x) \ge -1 \quad\text{and}\quad P_n'(x) \le 1 \qquad \forall x \in [0,n]. \tag{3.3}
\end{equation}

%% [Editorial note on a printed point (recorded, not corrected): (3.3) is printed in this
%%  non-strict form.  The strict negation of statement (1) would read "P_n(x) > -1 and
%%  P_n'(x) < 1"; the non-strict form printed here is what is reproduced, and it is what the
%%  authors pass to the limit with when they obtain (3.4).]

These assumptions easily imply that
\[
-1 \;\le\; P_n(x) \;\le\; P_n(0)+x \;=\; 1+x \;\le\; 2
\]

%% --- page 1010 (PDF page 10; running head: MARINA GHISI AND MASSIMO GOBBINO) ---

for every $x \in [0,1]$. In particular, the sequence $\{P_n(x)\}$ is bounded for at least $(k+1)$
distinct values of $x$, which we denote by $x_1,\dots,x_k,x_{k+1}$.

As a consequence, there exists a polynomial $P_\infty(x)$, with degree at most $k$, such that
$P_n(x) \to P_\infty(x)$ up to subsequences (not relabeled), in the sense that we have both
convergence of coefficients, and pointwise convergence. This is a well known fact which follows, for
example, from the simple remark that coefficients of a polynomial $P(x)$ with degree less than or
equal to $k$ depend in a linear way (through a Vandermonde matrix) from the values of $P(x)$ in the
$(k+1)$ points $x_1,\dots,x_k,x_{k+1}$. Alternatively, it follows from the fact that the maximum of
$|P(x_i)|$ for $i=1,\dots,k+1$, and the maximum of the absolute value of the coefficients of $P(x)$
are two norms on the space of all polynomials with degree less than or equal to $k$, and they are
equivalent because the space is finite dimensional.

In any case, convergence of coefficients implies that $P_\infty(0)=1$ and $P_\infty'(0)=-1$.
Pointwise convergence implies that we can pass to the limit in (3.3) and obtain that
\begin{equation}
P_\infty(x) \ge -1 \quad\text{and}\quad P_\infty'(x) \le 1 \qquad \forall x \ge 0. \tag{3.4}
\end{equation}
If the degree of $P_\infty(x)$ is $1$, then $P_\infty(x)=1-x$, which does not satisfy (3.4). If the
degree of $P_\infty(x)$ is at least $2$, then as $x \to +\infty$ we have that both $P_\infty(x)$ and
$P_\infty'(x)$ tend to $+\infty$ or $-\infty$ according to the sign of the leading coefficient. This
contradicts (3.4) also in this case, and completes the proof of statement (1).

Now we claim that the conclusion of statement (2) is true with
\[
\varepsilon_k := \min\left\{\frac{1}{\delta_k^{\,k+1}},\;\frac{1}{k\delta_k^{\,k}}\right\}.
\]
Indeed from statement (1) we know that there exists $x \in [0,\delta_k]$ such that either
$P(x) \le -1$ or $P'(x) \ge 1$. In the first case the first inequality in the left-hand side of (3.2)
gives that (clearly $\delta_k \ge 1$ for every $k \ge 1$)
\[
A \;\ge\; \frac{-P(x)}{x^{k+\alpha}} \;\ge\; \frac{1}{x^{k+\alpha}} \;\ge\; \frac{1}{\delta_k^{\,k+\alpha}}
\;\ge\; \frac{1}{\delta_k^{\,k+1}}.
\]
In the second case the second inequality in the left-hand side of (3.2) gives that
\[
A \;\ge\; \frac{P'(x)}{kx^{k+\alpha-1}} \;\ge\; \frac{1}{kx^{k+\alpha-1}} \;\ge\;
\frac{1}{k\delta_k^{\,k+\alpha-1}} \;\ge\; \frac{1}{k\delta_k^{\,k}}.
\]
In both cases implication (3.2) holds true. \hfill$\square$

Next result is the main step toward higher order Glaeser inequalities, both in $\mathbb{R}$ and in
bounded intervals. The proof exploits the ``magic constants'' of Lemma 3.3 and Taylor's expansion.

%% [Editorial note, not transcription: the paragraph just above is printed on p. 1010 as the
%%  opening of the next item (Proposition 3.4, printed p. 1011, already collected).  It is kept
%%  because it is printed on the page transcribed here and because it states, in the authors'
%%  own words, the dependency of Proposition 3.4 on Lemma 3.3.]

\end{document}
